\section{Related Work}
\label{sec:related}

% Draft for Ammaar's revision. Positioning rules (from the plan):
% - C5 must be positioned against ERC/convergence-aid prior art
%   honestly, NOT overclaimed as "unlike any software".
% - Claim "first PUBLISHED benchmark on Digitize-HCD", not "first to use".

\subsection{Hand-Drawn Circuit Recognition Pipelines}
Early work treated the task as symbol recognition plus ad-hoc wire
tracing. Reddy and Panicker \cite{reddy2022handdrawn} combine YOLOv5
component detection with Hough-based node recognition and report
98.2\% detection mAP@0.5, but evaluate connectivity only qualitatively
on a small private set. Recent end-to-end systems are stronger: SINA
\cite{sina2026} couples detection, connected-component connectivity
reasoning, OCR, and a vision-language model for designator assignment;
Netlistify \cite{netlistify} adds orientation detection to component
recognition. Wang et al.\ \cite{wang2026topology} is closest to our
setting: a topology-consistent parsing framework trained on
Digitize-HCD with wire connected-component-guided connectivity and
explicit endpoint semantics, reporting a 95.14\% strict image-level
success rate. In all of these, however, evaluation is performed on
private or unreleased test sets with bespoke metrics, so results
cannot be compared or audited. Our contribution is deliberately
complementary: a frozen, human-verified, public evaluation protocol
(C1/C4) plus baseline results, against which such systems can be
measured.

\subsection{Datasets and Graph Extraction on CGHD}
Two public datasets anchor the field. CGHD \cite{cghd} provides
hand-drawn circuit images by multiple drafters with symbol bounding
boxes, and later releases add segmentation maps; the accompanying line
of work by Bayer et al.\ extracts electrical graphs via instance
segmentation \cite{bayer2023instance}, modular pipelines
\cite{bayer2024modular}, and text extraction \cite{bayer2023text}.
Digitize-HCD \cite{digitizehcd_data} contributes 1{,}277 photographed
hand-drawn circuits with 18{,}600 component boxes over 17 classes,
text transcriptions, and — uniquely — per-class \emph{port-location}
annotations (Gaussian heatmaps and pixel coordinates). Neither dataset
ships connectivity ground truth for end-to-end evaluation; our
benchmark adds exactly that layer on Digitize-HCD, and we use CGHD's
class dictionary for a cross-dataset mapping. \draftnote{mention the
P\&ID/engineering-drawing digitization literature in one sentence if a
survey ref is added to refs.bib.}

\subsection{Making Recovered Circuits Simulate}
That a syntactically valid netlist can still fail DC analysis is a
classic problem: simulators ship convergence aids (gmin stepping,
source stepping, \texttt{.nodeset}) \cite{ngspice}, and schematic
tools ship electrical rule checks (ERC) that flag floating nodes or
missing references. Our C5 layer builds deliberately on this prior
art rather than claiming novelty over it: what is new is (i) applying
ERC-style diagnosis to \emph{recovered} (noisy, machine-inferred)
circuits rather than human-entered ones, (ii) enforcing a
\emph{minimality} discipline over the applied fixes, and (iii) the
\emph{assumption ledger} — a per-circuit, machine-readable record that
classifies each repair as behavior-invariant (\emph{gauge}) or
behavior-changing (\emph{assumption}), with alternatives, so a human
can audit exactly what the system assumed about their design intent.
To our knowledge no published hand-drawn-circuit system measures or
discloses this repair step; they either report simulation success
without stating what was fixed, or stop at topology.
